\scanpage{11}
$A'$ se réduisent bien, $A'$ se réduit bien. On
sait prouver la réciproque, qui \unsure{équivaut|s'exprime}
au

\uline{Théorème} (\unsure{Shimura} - Koizumi) Si $A'$ sur $K$
se réduit bien, tt schéma abélien
quotient aussi.
% slanted
\margin{On peut le prouver \uline{directement} comme la proposition 3}

On \deleted{peut \unsure{prouver}} suppose $A'=A\otimes_V K$,
et le quotient en question \uline{facteur} direct.
On a donc des endomorphismes
$p'\unsure{,|;}q'$ de $A'$ tels que
\[
p'\unsure{+}q'=\mathrm{id},\qquad p'{}^2=p',\qquad q'{}^2=p'
\]
Ces endomorphismes correspondent à
des endomorphismes $p\unsure{,|;}q$ de $A$, \unsure{satisfaisant|satisfont}
les $\unsure{\mathrm m^{es}}$ conditions.
\[
p+q=\mathrm{id},\qquad p^2=p,\quad q^2=q.
\]
Donc on aura un isom
\[
A=\mathrm{Ker}\,p\times\mathrm{Ker}\,q=A_1\times A_2
\]
% margin note written sideways
\margin{\deleted{\unsure{Condition|Conclusion|Corollaire}}}(\unsure{c'est} ``fonctoriel'' !). Donc $A_1$ et
$A_2$ sont propres sur $V$, d'ailleurs la
somme de leurs fibres en $k$ est $A\otimes k$,
qui est \deleted{\unsure{simple}}\inserted{abélien} sur $k$, il faut que $A_1\otimes k$ et $A_2\otimes k$
soient abéliens sur $k$, donc \unsure{aussi} $A_1$ et $A_2$. D'où \unsure{le|la}
% archive sheet number: 10
